% Crank-Nicolson stencil: one cell at t^n and t^{n+1} with neighbours.
\begin{tikzpicture}[
  font=\footnotesize,
  node distance=18mm,
  T/.style={circle, draw=black, line width=0.5pt, fill=white,
            minimum size=11mm, inner sep=0pt,
            font=\footnotesize},
  arr/.style={-{Stealth[length=1.8mm]}, line width=0.5pt, color=dim},
  arrCN/.style={-{Stealth[length=2mm]}, line width=0.8pt, color=teal},
  every node/.style={align=center}
]
% Time axis label
\node[font=\sffamily\bfseries, color=char] at (-30mm, 0)
  {time\\$t^n$};
\node[font=\sffamily\bfseries, color=char] at (-30mm, -20mm)
  {time\\$t^{n+1}$};
% t^n row
\node[T] (Tnm1) at (0,0) {$T_{i-1}^{n}$};
\node[T] (Tni)  at (20mm,0) {$T_{i}^{n}$};
\node[T] (Tnp1) at (40mm,0) {$T_{i+1}^{n}$};
% t^{n+1} row
\node[T] (Tnp1m1) at (0,-20mm)  {$T_{i-1}^{n+1}$};
\node[T] (Tnp1i)  at (20mm,-20mm) {$T_{i}^{n+1}$};
\node[T] (Tnp1p1) at (40mm,-20mm) {$T_{i+1}^{n+1}$};
% Coupling arrows (lateral, t^n)
\draw[arr] (Tnm1.east) -- (Tni.west);
\draw[arr] (Tnp1.west) -- (Tni.east);
% Coupling arrows (lateral, t^{n+1})
\draw[arr] (Tnp1m1.east) -- (Tnp1i.west);
\draw[arr] (Tnp1p1.west) -- (Tnp1i.east);
% Vertical CN time-coupling
\draw[arrCN] (Tni.south) -- (Tnp1i.north);
% Annotation: half + half on the right
\node[anchor=west, color=teal, font=\sffamily\small\bfseries]
  at (52mm, -10mm)
  {\begin{tabular}{@{}l@{}}CN \\ averages\\ the two\\ rows\end{tabular}};
\end{tikzpicture}
